\section{Related Work}
\label{sec:related-work}

Ocular-artifact correction has traditionally relied on an auxiliary EOG recording or on spatial separation within the EEG. Regression methods estimate how ocular potentials propagate to each scalp channel and subtract the predicted contribution \cite{gratton1983ocular,semlitsch1986ocular}. Frequency-domain regression and adaptive filtering extend this view to frequency-dependent or time-varying coupling, whereas independent component analysis and artifact subspace reconstruction separate artifacts from neural activity without an explicit EOG-to-EEG transfer model \cite{jung2000bss,chang2020asr}. Calibration-based subspace subtraction has shown that eye-related components can be estimated once and applied online through matrix operations \cite{kobler2020sgeyesub}. These methods establish the value of a calibration segment, but they do not combine participant-specific ocular propagation with a learned distribution of clean EEG.

Deep EEG denoisers learn nonlinear mappings from contaminated to clean or cleaner signals. EEGdenoiseNet introduced a widely used semi-simulated benchmark, and subsequent convolutional, decomposition-based, and U-Net models improved artifact removal under related paired constructions \cite{zhang2021eegdenoisenet,sun2020rescnn,yu2022deepseparator,chuang2022icunet,lai2024cleegn}. Their main strength is flexible waveform reconstruction; their usual limitation for the present setting is that participant variation is absorbed into shared model parameters. A participant label or a generic recording embedding does not directly specify the ocular component to remove. Our method instead retains a shared denoiser and conditions it on a propagation estimate obtained from simultaneous EEG and EOG during a separate calibration interval.

Diffusion models learn a reverse process from a prescribed noising process and can be conditioned on degraded observations for image restoration \cite{ho2020ddpm,song2021scoresde,saharia2022palette}. Methods such as DDRM and diffusion posterior sampling incorporate a known degradation process during reconstruction \cite{kawar2022ddrm,chung2023dps}; EEGDfus applies conditional diffusion directly to EEG denoising \cite{huang2025eegdfus}. Our formulation is closest in spirit to conditional restoration, but the condition is factored into a corrupted EEG observation and a separately calibrated EOG-derived artifact estimate. The method also returns sample-based predictive intervals. We evaluate these intervals through empirical marginal coverage and the continuous ranked probability score rather than treating stochastic variation as calibrated uncertainty by construction \cite{gneiting2007strictly}.
